\section{利得表は他者の行動ごとに比較する}
\label{theory:design:game}
各参加者が自分の利得を最大にするゲームでは、相手の行動を固定して自分の候補を比較する。相手のどの行動でも優る戦略が支配戦略である。両者の最適反応が同時に満たされる組合せは、片方だけ変更しても利得を増せないナッシュ均衡となる。行プレイヤーの利得を列ごと、列プレイヤーの利得を行ごとに比較し、表を転置して読む誤りを避ける。両者の合計利得の最大化は協調の評価で、独立の最適反応とは別である。

\section{期待値と決定ノードの後向き評価}
\label{theory:design:decision}
確率が既知の状態\(s\)について期待便益は\(\sum_s p_s B_s\)。一定費用\(C\)を引けば期待純便益になる。決定後に状態を観測して対応を選べる木では、末端から機会ノードを期待値、決定ノードを最大便益又は最小損失で置き換える。選べる二案を半々に平均してはならない。先に選ぶ案と、状態を見てから選べる案では利用できる情報が違う。

\section{限界利益と検査の期待費用}
\label{theory:design:economics}
一人又は一個当たり売上を\(p\)、変動費を\(v\)、期間固定費を\(F\)、目標利益を\(G\)とすると、必要数量は\((F+G)/(p-v)\)。\(p-v>0\)が前提である。最後に営業日数・設備数で割って、求める単位へ変換する。
\label{theory:design:inspection}
検査回数\(x\)で不具合確率が\(p(x)\)、一回費用\(c\)、損害額\(D\)なら期待総費用は\(cx+Dp(x)\)。検査を増やす費用と確率低下の利益を比較する。連続最適点を微分で求めた後、実際に許される整数値と端点を調べる。推定確率のモデルをそのまま現場の保証値とはしない。

\section{凸性は最適解の意味を強める}
\label{theory:design:convex}
線形計画は目的と制約がともに線形である。凸集合上で凸関数を最小化する凸計画では局所最適点が大域最適点となる\autocite{finish56_convex}。目的だけ凸でも実行可能集合が非凸ならその性質は保証されない。整数制約は実行可能集合を離散化し、連続問題の解を丸めるだけでは制約を満たさない場合がある。ヒューリスティックは候補解を得る手段で、最適性証明や誤差保証は別に扱う。

\section{経済的発注量は二つの費用を釣り合わせる}
\label{theory:design:eoq}
一定需要\(d\)、即時一括補充、欠品なしのモデルでは、発注量\(q\)に対し発注回数\(d/q\)、平均在庫\(q/2\)。一回費用\(k\)、年間単位在庫費\(h\)なら
\begin{equation}C(q)=\frac{kd}{q}+\frac{hq}{2},\qquad q_*=\sqrt{\frac{2kd}{h}}.\end{equation}
\(q>0\)で\(C''=2kd/q^3>0\)だから停留点は最小である\autocite{finish56_eoq}。上限、ロット単位、需要変動、納期、数量割引があれば再計算する。全需要に比例して一定の購入費は、このモデルでは最適発注量に影響しない。

\section{三つの解析公式の成立条件}
\label{theory:design:mathformulas}
Taylor展開の中心が0ならMaclaurin展開で\(f(x)=\sum_{k\ge0}f^{(k)}(0)x^k/k!\)。無限回微分可能だけでは級数と元の関数の一致は保証されず、剰余の収束を含む条件が必要である。複素指数の\(e^{i\theta}=\cos\theta+i\sin\theta\)から\(\theta=\pi\)で\(e^{i\pi}+1=0\)となる。
ロピタルの定理は\(0/0\)又は適切な無限大比等の不定形で、微分可能性、分母微分非零、微分比の極限等を確認して使う。任意の比の分子・分母を微分する規則ではない。

\section{不等式の向きと等号条件}
\label{theory:design:inequality}
非負の数の幾何平均は算術平均以下で、全数が同じなら等号を取る。\([0,\pi/2]\)上の\(\sin\theta\)の凹性から弦との比較で\(\sin\theta\ge2\theta/\pi\)。\(0<\theta\le\pi/2\)なら\(\sin\theta/\theta\ge2/\pi\)、右端で等号である。
一方、区間上で\(f''>0\)なら厳密凸で、異なる二点に対し\(f((x_1+x_2)/2)<(f(x_1)+f(x_2))/2\)。端点を含む不等式と厳密凸の厳密不等号を読み分ける。

\section{歯車の速度比とトルク}
\label{theory:design:gear}
回転動力は\(P=T\omega\)。速度比\(r=\omega_{\mathrm{in}}/\omega_{\mathrm{out}}\)、効率\(\eta\)なら\(T_{\mathrm{out}}=\eta rT_{\mathrm{in}}\)。減速で角速度が下がると、同じ動力に対応するトルクは増す。比の定義を入力/出力かその逆かで確認し、回転方向と大きさの比較を分ける。

\section{平面応力と平面ひずみは異なる近似}
\label{theory:design:plane}
薄板の自由表面に対する平面応力では\(\sigma_z,\tau_{xz},\tau_{yz}\simeq0\)で、面外ひずみは一般に0でない。軸方向変形を拘束した一様な長い構造の内部に対する平面ひずみでは\(\varepsilon_z,\gamma_{xz},\gamma_{yz}\simeq0\)で、軸応力は一般に残る。等方線形弾性なら後者の\(\sigma_z=\nu(\sigma_x+\sigma_y)\)。端部、軸方向の自由伸び、異方性を含む場合は境界条件からモデルを選び直す。
